\section{Discussion}
\label{sec:discussion}
% ------------------------------------------------------------
% Sumber: §13.8.3, §13.9, finding 07 §posisi-kita.
% ------------------------------------------------------------

\paragraph{Three properties, three literatures.}
``Can LLMs simulate populations?'' conflates three claims. Prior work
established property (i), \emph{readable} \citep{llmopinions2026}, and
documented the behavioral symptom from outside
(\citealp{sim2real2026,odyssim2026}). We establish property (ii),
\emph{faithfully arranged}---at specific locations, unevenly across
attribute types---and show that property (iii), \emph{causally used}, is
not a function of it. Under cluster-robust inference, causal use is not
absent---one type shows a clear, Bonferroni-surviving pathway at L11, and
two more are suggestive---but it does not follow the map
(Sec.~\ref{sec:sweep}). Natural-prompt simulation therefore does not fail
simply because the map is weak; a good map is no guarantee the model
consults it.

\paragraph{A mechanistic candidate for the Sim2Real gap---and a warning
about explaining it by representation quality.}
Behavioral homogeneity of LLM simulators is well measured from outside
\citep{sim2real2026}. Our results offer a consistent inside account:
outputs vary little across personas overall
(Sec.~\ref{sec:ceiling}), and forcing identity through the network moves
them only slightly under a deliberately strong instrument. What our
results do \emph{not} support is that simulation fails because the
internal representation is poor: homogeneity looks like at least two
mechanisms---insufficient causal throughput where the map is good, and no
good map to throughput where it is not. One qualification: the behavioral
literature evaluates post-trained assistants, whereas we measure a base
model; since post-training is itself a plausible cause of
homogenization, our results establish only that a representational
component of homogeneity predates post-training, rather than explaining
any particular assistant's behavior.

\paragraph{Why surface fine-tuning may not transfer.}
SubPOP-style fine-tuning \citep{subpop2025} improves in-distribution
accuracy but degrades on unseen surveys. If faithful structure already
exists internally, output-surface fine-tuning may re-teach what the model
already encodes---a survey-specific surface mapping instead of a
connection to the existing map. This is a testable hypothesis, not a
demonstrated mechanism; the direct test (does reading from the faithful
location transfer where surface fine-tuning does not?) is future work.

\paragraph{Practical guidance now.}
For population simulation: check second-order fidelity before relating
groups, treat demographic attributes as heterogeneous
(Sec.~\ref{sec:act2}), and measure identity conditioning rather than
assume it. \textbf{Do not steer through a single apparently-faithful unit
without a corruption control}: our causal analysis's most consistent
effect is a \emph{harmful} one from doing exactly this
(RACE$\times$RELIG, \starhead{} alone, wrong direction,
$p_{\text{pair}}=1/4096$; Sec.~\ref{sec:v2}). Read a per-group
distribution rather than ask for it: a linear probe on \starhead{} alone
lands $21$--$31\%$ closer to survey truth than the model's own answers
(Sec.~\ref{sec:probe}), yet a group-blind baseline still wins on
per-question ordering.
